%!TEX root = ../main.tex

\section{Conclusion}\label{sec:conclusion}

We develop a structural theory of endogenous premium formation for Bitcoin
treasury vehicles.  The NAV premium is the market value of a perpetual
accumulation option---the right to issue equity above NAV and purchase BTC
accretively.  The fair premium satisfies a fixed-point equation whose
steady-state closed form connects to the Black--Scholes formula, linking
the premium to BTC volatility, issuance capability, leverage, and discount
rates.

The reflexive feedback between premium and accumulation generates multiple
equilibria: a high-premium flywheel, an unstable tipping point, and a
low-premium death spiral, connected by saddle-node bifurcations with
hysteresis.  We reinterpret five structural indicators as the option Greeks
of the accumulation claim and derive optimal capital structure rules
governing the transition among financing channels.

Calibrating to Strategy (MSTR) through March~2026, we find the current
premium of 24.2\% sits below the OU long-run mean ($\text{PMRI} = -0.71$),
implying positive carry and an underpriced accumulation option.  Leverage is
moderate ($\mathrm{ILE} = 1.49$), the flywheel is active
($\mathrm{IBGR} = 17.6\%$), and three-year survival probability is 91.9\%.
The framework provides a unified lens for valuation, risk monitoring, and
vehicle selection in the emerging landscape of corporate Bitcoin treasury
strategies.
